%=======================================================================
%  CHAPTER 6 — SPEECH AND CHANNEL CODING FUNDAMENTALS  (4 hrs)
%=======================================================================
\section{Speech and Channel Coding Fundamentals}

{\color{sub}\footnotesize 4 hrs \gap$\cdot$\gap asked in \mk{22 of 22 papers}
\gap$\cdot$\gap 5 syllabus sub-topics \gap$\cdot$\gap a few worked examples
$\rightarrow$ Ch 6 numerical companion.}

\lead{Read this first:}
\begin{itemize}
  \item \textbf{Speech coding is the single most reliable topic in the whole paper.} Some part of
        it appears in \mk{21 of 22 papers}, and \textbf{characteristics of the speech signal}
        \tS{7} plus \textbf{vocoders} \tS{10} plus \textbf{LPC} \tS{8} are the three that recur.
  \item Almost every speech question is a \textbf{pair}: characteristics of speech
        \textbf{plus} one coder. So learn the ACF / PDF / PSD trio once, then attach whichever
        coder is asked.
  \item \textbf{Channel coding} is examined much more lightly, in \tF{9} papers but never for more
        than \m{5}. Viterbi \tF{4} and turbo \m{5} are the recurring ones.
  \item \mk{The one link to remember between the two halves:} speech coding \textbf{removes}
        redundancy to save bandwidth, and channel coding then \textbf{adds} redundancy back to
        survive the channel. They pull in opposite directions on purpose.
\end{itemize}

\hr

\T{6.1 Speech coding: why, and the two families}

\Q{Why do we need speech coding techniques?}{\m{4}\m{2}}
{\color{sub}\footnotesize \yr{70 Ma, \textbf{78 Ch}} \gap$\cdot$\gap
the framing paragraph for every other question in this section}

\sbsr{0.55}{%
\lead{Goals of speech coding}
\begin{itemize}
  \item \textbf{Remove redundancy} from the speech signal.
  \item \textbf{Achieve compression}: use the minimum number of bits while maintaining quality.
  \item \textbf{Reduce transmission and storage cost}.
  \item In wireless specifically: \mk{transmit speech at the highest possible quality using the
        least possible channel capacity}, so more speech channels fit in a given bandwidth.
\end{itemize}

\lead{The trade-offs}
\begin{itemize}
  \item Higher speech bit rate $\rightarrow$ higher quality, but greater bandwidth and storage.
  \item So a balance is struck between \textbf{bit-rate efficiency} and
        \textbf{algorithmic complexity}.
  \item A third axis matters in mobile: \textbf{robustness to channel errors}. A very low rate
        coder throws away so much that a single bit error is audible.
\end{itemize}

\lead{The two families \mk{(state this split before naming any coder)}}
\begin{itemize}
  \item \textbf{Waveform coders}
  \begin{itemize}
    \item \textbf{No speech model is required.} Work on \textbf{any} input analog signal.
    \item Designed to be \textbf{source independent}, so they code music and tones equally well.
    \item Aim: reproduce the \textbf{time waveform} as closely as possible, minimising the bits
          per sample.
    \item \textbf{($+$)} Robust across a wide range of speakers and noisy environments.
          \textbf{Minimal complexity.}
    \item \textbf{($-$)} Only \textbf{moderate} bit-rate economy, ie \mk{low compression ratio}.
    \item Eg \textbf{PCM, DPCM, ADPCM, delta modulation, CVSDM}, and the frequency-domain
          coders below.
  \end{itemize}
  \item \textbf{Source coders (vocoders)}
  \begin{itemize}
    \item Use \textbf{a priori knowledge} of how speech is produced.
    \item Recreate the \textbf{spectral characteristics} of speech by modelling the vocal tract
          as a \textbf{time-varying digital filter}.
    \item Filter parameters are extracted at the encoder, \textbf{only the parameters are sent},
          and the decoder \textbf{re-synthesises} speech from them.
    \item \textbf{($+$) Very high economy in bit rate}, down to 2.4 kbps.
    \item \textbf{($-$) High complexity}, and quality is low and \mk{mechanical sounding},
          though intelligible. Performance is also \textbf{talker dependent}.
  \end{itemize}
  \item \textbf{Hybrid coders} sit between the two and are what modern cellular actually uses.
        \added
\end{itemize}}{%
\figT{c6_hierarchy.png}
\vspace{1pt}
{\color{sub}\footnotesize The hierarchy of speech coders. \mk{Draw this tree and the ``types of
speech coder'' mark is banked.}}}

\Q{Explain the characteristics of the speech signal, and how they are used in designing coders}{\m{8}\m{6}\m{4}\m{3}}
{\color{sub}\footnotesize \tS{7} \yr{\textbf{72 Ash}, \textbf{77 Ch}, 73 Ma, \texttt{79 Bh},
\textbf{80 Ch}, \textbf{\texttt{82 Bh}}, \texttt{81 Ba}} \gap$\cdot$\gap
\mk{the most repeated question in the chapter, and it pairs with almost every other one}}

{\color{sub}\footnotesize Speech waveforms have properties that can be \textbf{exploited} when
designing efficient coders. Three are used most often. For each one, always say
\textbf{what the property is} and then \textbf{which coder design it justifies}: that second half
is where the marks are.}

\vspace{2pt}
\begin{tabularx}{\textwidth}{@{}L{3.5cm}XX@{}}
\toprule
\hrow \thead{Property} & \thead{What it says} & \thead{What it justifies} \\
\midrule
\textbf{Autocorrelation function (ACF)} &
  Adjacent samples of a speech segment are \textbf{strongly correlated}. Typical adjacent-sample
  correlation is \mk{0.85 to 0.9}. So every sample is largely \textbf{predictable} from previous
  samples, plus a small random error.
  \[ C(k) = \frac{1}{N}\sum_{n=0}^{N-|k|-1} x(n)\,x(n+|k|) \] &
  \textbf{All differential and predictive coding schemes}: DPCM, ADPCM, delta modulation, and
  \textbf{LPC}. Send the small prediction error instead of the sample. \\
\midrule
\textbf{Probability density function (PDF)} &
  Speech amplitude has a \textbf{highly non-uniform} pdf: very high probability of
  \textbf{near-zero} amplitudes, small but significant probability of very high amplitudes, and a
  \textbf{monotonically decreasing} function in between.
  \[ p(x) = \frac{1}{\sqrt{2}\,\sigma_x}\,e^{-\sqrt{2}\,|x|/\sigma_x} \] &
  \textbf{Non-uniform quantizers} and \textbf{vector quantizers}: put \mk{more quantization levels
  where the probability is high} and fewer where it is low. This is why $\mu$-law and A-law
  companding exist. \\
\midrule
\textbf{Power spectral density (PSD)} &
  The PSD of speech is strongly \textbf{non-flat}: energy is concentrated at low frequencies.
  High-frequency components carry \textbf{little energy but a lot of information}. &
  \textbf{Frequency-domain coding}: code different frequency bands separately and allocate bits
  per band. \mk{Significant coding gain}, but high bands must still be adequately represented or
  intelligibility collapses. \\
\bottomrule
\end{tabularx}

\vspace{3pt}
\sbs{%
\lead{Two more properties worth a line}
\begin{itemize}
  \item \textbf{Band-limited}: speech is essentially confined to \textbf{300 Hz to 3400 Hz},
        so 8 kHz sampling suffices. This is what makes digital speech possible at all. \added
  \item \textbf{Quasi-stationary}: speech statistics are roughly constant over
        \textbf{10 to 30 ms}, which is why every coder works on frames of that length. \added
\end{itemize}}{%
\lead{Voiced Vs unvoiced speech +Fig}
\begin{itemize}
  \item \textbf{Voiced} sounds (``m'', ``v'', all vowels) are produced by
        \textbf{quasi-periodic} vocal cord vibration. High energy, clear \textbf{pitch period}.
  \item \textbf{Unvoiced} sounds (``s'', ``f'') are produced by turbulent air, so they look like
        \textbf{random noise}. Low energy, no pitch.
  \item \mk{Every vocoder must transmit a voiced / unvoiced decision plus the pitch period.}
\end{itemize}
\vspace{2pt}
\figT{c6_speech_wave.png}}

\hr

\T{6.2 Frequency domain coding of speech}

\Q{Explain the types of frequency domain coding of speech. Explain sub-band coding}{\m{4}\m{6}}
{\color{sub}\footnotesize \tP{2} \yr{\textbf{73 Bh}, \textbf{\texttt{80 Bh}}} \gap$\cdot$\gap
80 Bh asks for sub-band coding \mk{with transmitter and receiver}}

\sbsr{0.52}{%
\lead{The idea}
\begin{itemize}
  \item Frequency-domain coding is a \textbf{waveform} coding method.
  \item Think of it as \mk{distributing quantization noise across the signal spectrum} on purpose,
        rather than letting it spread uniformly.
  \item The speech signal is split into a set of \textbf{frequency components}, which are
        \textbf{quantized and encoded separately}.
  \item The bits allocated to each component can be \textbf{varied dynamically} and shared between
        bands.
  \item Two algorithms are examined: \textbf{sub-band coding} and \textbf{adaptive transform
        coding}.
\end{itemize}

\lead{Sub-band coding (SBC)}
\begin{itemize}
  \item Divides speech into \textbf{4 or 8 sub-bands} using a \textbf{bank of filters}.
  \item Each sub-band is \textbf{sampled at its bandpass Nyquist rate}, then encoded
        \textbf{separately}.
  \item Aim: \textbf{control and distribute the quantization noise} over the whole spectrum.
  \begin{itemize}
    \item Noise generated in one band \textbf{stays in that band}.
    \item So a low-energy band is not swamped by noise leaking from a high-energy one.
  \end{itemize}
  \item Bits per band are allocated by \textbf{perceptual importance}: low bands carry most of the
        energy and get the most bits.
  \item \textbf{Bit rate range: 9.6 kbps to 32 kbps.}
  \item \textbf{Implementation trick:} each sub-band is \textbf{low-pass translated to zero
        frequency} using a scheme like single-sideband modulation. That
        \mk{reduces the required sampling rate} for each band.
  \item \textbf{Receiver}: decode each band, \textbf{interpolate} back up, band-pass translate to
        the original position, and \textbf{sum}.
  \item \textbf{($+$)} Good quality at moderate rates, and the split is easy to implement.
        \textbf{($-$)} Filter bank adds delay and complexity.
\end{itemize}}{%
\figT{c6_subband.png}
\vspace{1pt}
{\color{sub}\footnotesize Sub-band coder and decoder. LP translator, decimator and encoder per
band on the left, the inverse chain on the right.}

\lead{Adaptive transform coding (ATC) \pill{gdL}{gd}{syllabus, no PYQ yet}}
\begin{itemize}
  \item Applies a \textbf{block transform} to windowed segments of the speech waveform.
  \item Each segment becomes a set of \textbf{transform coefficients}, quantized separately with a
        number of bits \textbf{proportional to its perceptual significance}.
  \item \textbf{Bit rate range: 9.6 kbps to 20 kbps.}
  \item Most common transform: the \textbf{discrete cosine transform (DCT)}
        \[ \begin{gathered} X_c(k) = \sum_{n=0}^{N-1} x(n)\,g(k)\,
           \cos\!\left[\frac{(2n+1)k\pi}{2N}\right], \\
           k = 0,1,\dots,N-1 \end{gathered} \]
        \[ x(n) = \frac{1}{N}\sum_{k=0}^{N-1} X_c(k)\,g(k)\,
           \cos\!\left[\frac{(2n+1)k\pi}{2N}\right] \]
        with $g(0)=1$ and $g(k)=\sqrt{2}$ for $k = 1,\dots,N-1$.
  \item \mk{SBC splits with filters, ATC splits with a transform.} That one line answers
        ``compare them''.
\end{itemize}}

\hr

\T{6.3 Vocoders}

\Q{What is a vocoder? Explain its operation with a block diagram, and the different kinds}{\m{2+6}\m{2+3}\m{3+5}\m{4}\m{6}}
{\color{sub}\footnotesize \tS{10} \yr{\textbf{70 Bh}, 70 Ma, \textbf{71 Bh}, \textbf{73 Bh},
73 Ma, \textbf{74 Bh}, 74 Ma, \textbf{78 Ch}, \textbf{\texttt{81 Bh}}, \texttt{81 Ba}}
\gap$\cdot$\gap \mk{the single most repeated question in the whole chapter}}

\sbsr{0.52}{%
\lead{What a vocoder does, in three steps}
\begin{itemize}
  \item \textbf{Analyze} the voice signal at the transmitter.
  \item \textbf{Transmit only the parameters} derived from that analysis, never the waveform.
  \item \textbf{Synthesize} the voice at the receiver from those parameters.
\end{itemize}
\begin{itemize}
  \item All vocoders model \textbf{speech generation as a dynamic system} and quantify the
        \textbf{physical constraints} of that system.
  \item Those constraints then give a very compact description of the signal.
  \item \textbf{($+$) Very high economy in transmission bit rate.}
  \item \textbf{($-$)} Much more complex than a waveform coder, \textbf{less robust}, and
        performance is \mk{talker dependent}.
\end{itemize}

\lead{The speech generation model +Fig \mk{(draw this first, always)}}
\begin{itemize}
  \item Sound is \textbf{generated} by an excitation source, then \textbf{modulated} as it passes
        through the vocal tract.
  \item So the synthesiser needs exactly two things:
  \begin{itemize}
    \item a sequence of \textbf{excitation signals}, and
    \item a \textbf{vocal tract approximation} to modulate them.
  \end{itemize}
  \item \textbf{Excitation} is switched by the voiced / unvoiced decision:
  \begin{itemize}
    \item \textbf{voiced} $\rightarrow$ a \textbf{pulse source}, a series of periodic pulses at
          the pitch rate,
    \item \textbf{unvoiced} $\rightarrow$ a \textbf{noise source}, a random sequence.
  \end{itemize}
  \item \textbf{Vocal tract} is modelled as a \textbf{linear filter}.
  \item Parameters transmitted: \textbf{pitch period}, \textbf{voiced / unvoiced flag},
        \textbf{pole frequencies} of the filter, \textbf{amplitude / gain}.
  \item \mk{Every vocoder below is a different way of describing that filter.} Say that and the
        four types fall out naturally.
\end{itemize}}{%
\figT{c6_speech_gen.png}
\vspace{1pt}
{\color{sub}\footnotesize The traditional speech generation model, the basis of \emph{all}
vocoding systems. Noise source and pulse source feed a vocal tract filter.}

\lead{The four classical vocoders \mk{(plus LPC)}}
\begin{itemize}
  \item \textbf{Channel vocoder} $=$ describes the filter by its \textbf{spectral envelope}, band
        by band.
  \item \textbf{Formant vocoder} $=$ describes it by the \textbf{positions of the spectral peaks}.
  \item \textbf{Cepstrum vocoder} $=$ separates source from filter by
        \textbf{homomorphic (log-spectrum) processing}.
  \item \textbf{Voice-excited vocoder} $=$ avoids describing the excitation at all, by sending the
        low band as PCM.
  \item \textbf{Linear predictive coder (LPC)} $=$ describes the filter by
        \textbf{prediction coefficients} in the \textbf{time} domain.
        \mk{The most popular of all.}
\end{itemize}}

\creamq{Channel vocoder}{\m{4}}
\sbsr{0.50}{%
\begin{itemize}
  \item The \textbf{first} analysis-synthesis system ever built.
  \item A \textbf{frequency domain} vocoder: it determines the \textbf{envelope} of the speech
        signal for a number of frequency bands, then samples, encodes and multiplexes those
        samples.
  \item Also transmits, per frame, the \textbf{voiced / unvoiced decision} and the
        \textbf{pitch frequency}.
  \item Sampling is done every \textbf{10 ms to 30 ms}.
\end{itemize}

\lead{Analyzer}
\begin{itemize}
  \item A bank of \textbf{bandpass filters}, each \textbf{100 Hz to 300 Hz} wide.
  \item Each filter output is \textbf{rectified} then \textbf{low-pass filtered}
        $\rightarrow$ this pair performs \textbf{envelope detection}.
  \item The low-pass bandwidth is chosen to \mk{match the rate at which the vocal tract changes},
        typically a few tens of Hz.
  \item Each envelope is \textbf{A/D converted} and fed to the encoder.
  \item In parallel, a \textbf{voicing detector} and a \textbf{pitch estimator} run on the input.
\end{itemize}

\lead{Synthesizer}
\begin{itemize}
  \item Decoder splits the stream back into band envelopes plus voicing and pitch.
  \item \textbf{Switch} selects the excitation:
        \textbf{pulse generator} driven by the received pitch period for voiced,
        \textbf{random noise generator} for unvoiced.
  \item Excitation is \textbf{multiplied} by each band envelope, then passed through the matching
        \textbf{bandpass filter}.
  \item All bands are \textbf{summed} to give the output speech.
\end{itemize}}{%
\figT{c6_chvoc_ana.png}
\vspace{1pt}
{\color{sub}\footnotesize Channel vocoder \textbf{analyzer}. Bandpass filter $\rightarrow$
rectifier $\rightarrow$ lowpass filter is the envelope detector.}
\vspace{4pt}
\figT{c6_chvoc_syn.png}
\vspace{1pt}
{\color{sub}\footnotesize Channel vocoder \textbf{synthesizer}. The switch picks pulse or noise
excitation from the received voicing information.}}

\creamq{Formant vocoder}{\m{4+4}}
{\color{sub}\footnotesize \tP{1} \yr{\textbf{71 Bh}} \gap$\cdot$\gap asked as
``formant vocoder operation \textbf{plus} characteristics of speech''}

\sbsr{0.50}{%
\begin{itemize}
  \item \textbf{Formants} $=$ the \textbf{spectral peaks} of the sound spectrum $|P(f)|$.
  \begin{itemize}
    \item They are frequency peaks with a \textbf{high degree of energy}.
    \item Especially prominent in \textbf{vowels}.
    \item Each corresponds to a \textbf{resonance of the vocal tract}.
    \item Roughly \mk{one formant every 1000 Hz}.
  \end{itemize}
  \item Can be viewed as a \textbf{type of channel vocoder} that estimates only the
        \textbf{first three or four formants} in a segment of speech.
  \item \mk{Key difference from the channel vocoder}: instead of sending samples of the whole
        power spectrum envelope, it sends only the \textbf{positions of the peaks}.
  \item It must also control the \textbf{intensities} of those formants, and send the
        \textbf{pitch period} and \textbf{voicing}.
  \item Per formant $k$ it transmits $F_k$ (the frequency) and $B_k$ (the bandwidth).
  \item \textbf{($+$)} Far fewer parameters than a channel vocoder $\rightarrow$
        \mk{much lower bit rate}.
  \item \textbf{($-$)} Reliable \textbf{formant tracking is hard}, especially for unvoiced speech
        and for nasals, so quality suffers.
\end{itemize}}{%
\figT{c6_formants.png}
\vspace{1pt}
{\color{sub}\footnotesize Spectral envelope of an [i] spoken by a male speaker. $F_1$, $F_2$ and
$F_3$ are the first three formants.}
\vspace{3pt}
\figT{c6_formant_ana.png}
\vspace{1pt}
{\color{sub}\footnotesize Formant vocoder \textbf{analyzer}: extracts $F_1..F_3$, their bandwidths
$B_1..B_3$, plus the voiced/unvoiced flag and pitch $F_0$.}
\vspace{3pt}
\figT{c6_formant_syn.png}
\vspace{1pt}
{\color{sub}\footnotesize Formant vocoder \textbf{synthesizer}: the excitation signal is shaped by
three resonators and summed.}}

\sbs{%
\lead{Cepstrum vocoder}
\begin{itemize}
  \item Speech is a \textbf{convolution} of an excitation source and the vocal tract system.
        Separating them normally needs deconvolution.
  \item \textbf{Objective of cepstral analysis:} separate the two \mk{without any prior knowledge
        of either}.
  \item \textbf{Cepstrum} $=$ the \textbf{inverse Fourier transform of the logarithm} of the
        estimated spectrum.
  \item The trick is that the log turns a product into a sum:
        \[ X(\omega) = E(\omega)\,V(\omega)
           \;\Longrightarrow\;
           \log X(\omega) = \log E(\omega) + \log V(\omega) \]
  \item After the inverse transform the two land in \textbf{different regions}:
  \begin{itemize}
    \item \textbf{low-frequency (low-quefrency)} coefficients $\rightarrow$ the
          \textbf{vocal tract spectral envelope},
    \item \textbf{high-frequency} coefficients $\rightarrow$ the \textbf{excitation}, a periodic
          pulse train at multiples of the pitch period.
  \end{itemize}
  \item So they can simply be \textbf{separated by windowing} (liftering).
  \item \textbf{Transmit only the vocal tract cepstral coefficients.}
  \item At the receiver: Fourier transform them to get the \textbf{vocal tract impulse response},
        then \textbf{convolve} that with a synthetic excitation to reconstruct the speech.
\end{itemize}}{%
\lead{Voice-excited vocoder}
\begin{itemize}
  \item \textbf{Hybrid}: \textbf{PCM} transmission of the \textbf{lower} frequency band of speech,
        combined with \textbf{channel vocoding} of the \textbf{higher} bands.
  \item At the synthesizer a \textbf{pitch signal is generated} from the received baseband by
        \textbf{rectifying, bandpass filtering and clipping} it. That creates a
        \textbf{spectrally flattened} signal with energy at the pitch harmonics.
  \item \mk{($+$) This eliminates the need for pitch extraction and voicing detection}, which are
        the two most error-prone operations in any vocoder.
  \item \textbf{($-$)} Costs bit rate: designed for \textbf{7.2 kbps to 9.6 kbps}, well above a
        pure vocoder.
\end{itemize}

\lead{Comparing the four}
\vspace{2pt}
\begin{tabularx}{\linewidth}{@{}L{1.95cm}XL{1.6cm}@{}}
\toprule
\hrow \thead{Vocoder} & \thead{What it sends} & \thead{Domain} \\
\midrule
Channel & Envelope per band, pitch, voicing & Frequency \\
Formant & 3--4 formant frequencies + bandwidths, pitch, voicing & Frequency \\
Cepstrum & Low-quefrency cepstral coefficients & Cepstral \\
Voice-excited & Baseband as PCM + high bands as channel vocoder & Hybrid \\
LPC & Prediction coefficients, gain, pitch, voicing & \textbf{Time} \\
\bottomrule
\end{tabularx}}

\hr

\T{6.4 Linear predictive coders}

\Q{Explain the linear predictive coder (LPC) with a neat block diagram}{\m{2+6}\m{3+5}\m{6}}
{\color{sub}\footnotesize \tS{8} \yr{\texttt{80 Ba}, \texttt{81 Ba}, 72 Ma, \texttt{79 Bh},
73 Ma, \textbf{76 Bh}, \textbf{75 Bh}, 71 Ma} \gap$\cdot$\gap
\mk{almost always paired with ``characteristics of the speech signal''}}

\sbsr{0.52}{%
\lead{What makes LPC different}
\begin{itemize}
  \item The \textbf{time domain} class of vocoders. All the others above work in frequency.
  \item Extracts the significant features of speech \textbf{from the time waveform} directly.
  \item \textbf{Computationally intensive}, but \mk{by far the most popular of the low bit rate
        vocoders}.
  \item Transmits \textbf{good quality voice at 4.8 kbps}, and intelligible speech at 2.4 kbps.
\end{itemize}

\lead{The model}
\begin{itemize}
  \item Models the vocal tract as an \textbf{all-pole linear filter}:
        \[ H(z) \;=\; \frac{G}{1 + \displaystyle\sum_{k=1}^{M} b_k z^{-k}} \]
        $G$ $=$ filter gain, $b_k$ $=$ filter (prediction) coefficients, $M$ $=$ order,
        typically 10.
  \item \textbf{All-pole is enough} because the vocal tract is a tube with resonances, and
        resonances are poles. \mk{That is the physical justification the examiner wants.}
  \item Coefficients are found by \textbf{minimum mean square error (MMSE)}:
        minimise the average energy in the \textbf{error signal}, the difference between
        \textbf{predicted} and \textbf{actual} speech amplitude.
  \item That works because of the \textbf{ACF property}: adjacent speech samples are 0.85 to 0.9
        correlated, so a short-order predictor already explains most of the signal.
\end{itemize}

\lead{Operation}
\begin{itemize}
  \item \textbf{Encoder}: buffer a frame $\rightarrow$ run \textbf{LP analysis} for the
        coefficients $\rightarrow$ decide \textbf{voiced / unvoiced} $\rightarrow$ estimate
        \textbf{pitch} $\rightarrow$ compute \textbf{gain}.
  \item Speech type, pitch, filter coefficients and gain are \textbf{coded into a bit stream} and
        transmitted.
  \item \textbf{Decoder}: pick pulse or noise excitation, scale by the gain, drive $H(z)$
        $\rightarrow$ \textbf{synthesized speech}.
\end{itemize}}{%
\figT{c6_lpc.png}
\vspace{1pt}
{\color{sub}\footnotesize Block diagram of an LPC coding system. Three parallel paths: LPC filter
coefficients, voiced/unvoiced decision, pitch. The synthesizer recombines them.}

\lead{Types of LPC \mk{(71 Ma asks for the types, \m{6})}}
\begin{itemize}
  \item Plain LPC uses a \textbf{single pulse per pitch period}, which sounds
        \mk{buzzy and synthetic}. The three variants below all improve the
        \textbf{excitation model}, and that is the only thing that separates them.
  \item \textbf{Multipulse excited LPC (MPE-LPC)}: uses \textbf{several pulses per frame}, with
        positions and amplitudes chosen by analysis-by-synthesis. No voiced / unvoiced decision is
        needed at all.
  \item \textbf{Residual excited LPC (RELP)}: see below.
  \item \textbf{Code excited LPC (CELP)}: the excitation sequence is selected from a
        \textbf{codebook of zero-mean Gaussian sequences}, shared by encoder and decoder.
  \begin{itemize}
    \item \textbf{Analysis by synthesis}: the encoder tries every codebook entry, synthesises
          speech with it, and keeps the one giving the smallest perceptually weighted error.
    \item Only the \textbf{codebook index} is transmitted, which is why it is so efficient.
    \item Bit rate \textbf{4800 bps}.
  \end{itemize}
\end{itemize}}

\sbsr{0.52}{%
\lead{Residual excited LPC (RELP)}
\begin{itemize}
  \item Improves quality \textbf{at the expense of a higher bit rate}, by computing and
        transmitting a \textbf{residual error}, exactly as DPCM does.
  \item Steps:
  \begin{itemize}
    \item LPC model and excitation parameters are estimated from a \textbf{frame of speech}.
    \item Speech is \textbf{synthesized at the transmitter} using them.
    \item The synthesized speech is \textbf{subtracted from the original} to form the
          \textbf{residual error}.
    \item The residual is \textbf{quantized, coded and transmitted} alongside the parameters.
  \end{itemize}
  \item At the receiver: synthesize from the model, then \textbf{add the residual back}.
  \item \mk{The receiver therefore gets the part of the signal the model could not explain},
        which is what lifts the quality above plain LPC.
\end{itemize}}{%
\figT{c6_relp.png}
\vspace{1pt}
{\color{sub}\footnotesize RELP encoder. The residual is formed by subtracting locally synthesized
speech from the buffered input.}}

\creamq{The GSM codec}{\m{2}}
{\color{sub}\footnotesize \yr{74 Ma} \gap$\cdot$\gap asked as part of a
``waveform + voice coding + GSM CODEC +Fig'' question}

\sbsr{0.50}{%
\begin{itemize}
  \item GSM uses the \textbf{regular pulse excited, long term prediction (RPE-LTP)} codec.
  \item It \textbf{combines} the earlier French \textbf{baseband RELP} codec with the German
        \textbf{multipulse excited long-term prediction (MPE-LTP)} codec.
  \item \mk{So it is a hybrid}: a source model (LPC) driving a partly waveform-coded excitation.
  \item Chain, left to right:
  \begin{itemize}
    \item \textbf{Pre-processing}: Hamming window, segmentation into 20 ms frames, pre-emphasis.
    \item \textbf{STP} (short term prediction): LPC inverse filter and STP analysis filter produce
          the residual $e_n$ and the \textbf{LAR($k$)} log-area-ratio coefficients.
    \item \textbf{LTP} (long term prediction): removes the \textbf{pitch periodicity} left in the
          residual, giving lag $P_n$ and gain $g_n$.
    \item \textbf{RPE}: low-pass filter, then \textbf{grid selection} picks the best regular pulse
          sequence $x_m$.
    \item All parameters are \textbf{multiplexed} into the coded speech stream.
  \end{itemize}
  \item Output rate \textbf{13 kbps} (full rate GSM), from 104 kbps of 13-bit input samples.
        \added
  \item \textbf{($-$) Relatively complex and power hungry.}
\end{itemize}}{%
\figT{c6_gsm_codec.png}
\vspace{1pt}
{\color{sub}\footnotesize GSM RPE-LTP speech encoder: pre-processing, short term prediction, long
term prediction, regular pulse excitation, multiplexer.}}

\hr

\T{6.5 Channel coding}

\Q{What is channel coding? Explain block codes and the definitions}{\m{2}}
{\color{sub}\footnotesize \yr{71 Ma} \gap$\cdot$\gap the definitions below are needed by every
other question in this section}

\sbsr{0.53}{%
\lead{The basic idea}
\begin{itemize}
  \item Add a \textbf{controlled amount of redundancy} to the message \textbf{before} it goes
        through the noisy channel.
  \item The redundancy is extra symbols added \textbf{in a known manner}, so the receiver can use
        them to \textbf{detect} and often \textbf{correct} errors.
  \item \textbf{Everyday analogy}: English already carries redundancy, which is why
        \mk{``CODNG THEORY IS AN INTRSTNG SUBJCT'' is still readable}.
  \item \mk{Note the tension with Ch 6.1}: speech coding strips redundancy out, channel coding
        puts a different, engineered redundancy back in.
  \item Improves \textbf{small-scale link performance}: faded data can still be recovered.
\end{itemize}

\lead{Definitions \mk{(learn these five)}}
\begin{itemize}
  \item A \textbf{word} is a sequence of symbols. A \textbf{code} is a set of vectors called
        \textbf{codewords}.
  \item \textbf{Hamming weight} $w(c)$ $=$ number of \textbf{non-zero elements} in a codeword.
        Eg $w(10110) = 3$.
  \item \textbf{Hamming distance} $d(c_1,c_2)$ $=$ number of places the two codewords
        \textbf{differ}. Eg $d(10110, 11011) = 3$.
        Note $d(c_1,c_2) = w(c_1 - c_2)$.
  \item \textbf{Minimum distance} $d'$ $=$ the smallest Hamming distance between any two
        codewords, $d' = \min_{i \ne j} d(c_i,c_j)$.
  \item \textbf{Code rate} of an $(n,k)$ code $= k/n$. It is the fraction of the codeword that is
        actual information.
\end{itemize}

\lead{Error detection and correction capability}
\begin{itemize}
  \item To \textbf{detect} $t$ errors per block: \; $d' \ge t+1$
  \item To \textbf{correct} $t$ errors per block: \; $d' \ge 2t+1$
  \item \mk{So minimum distance is the single number that decides how strong a code is.}
\end{itemize}}{%
\lead{Block codes}
\begin{itemize}
  \item A \textbf{block code} is a set of \textbf{fixed length} codewords. That fixed length is
        the \textbf{block length} $n$.
  \item A block code of size $M$ over a $q$-symbol alphabet is a set of $M$ $q$-ary sequences of
        length $n$.
  \item If $q = 2$ the symbols are \textbf{bits} and the code is \textbf{binary}.
  \item Usually $M = q^k$, and we call it an \textbf{$(n,k)$ code}: $k$ information symbols in,
        $n$ out, so \textbf{$n-k$ redundant symbols}.
  \item Block codes are \textbf{forward error correction (FEC)} codes: no return channel is needed.
\end{itemize}

\lead{Linear block codes (LBC)}
\begin{itemize}
  \item A code is \textbf{linear} if:
  \begin{itemize}
    \item the \textbf{sum of any two codewords} is also a codeword,
    \item the \textbf{all-zero codeword} is always a codeword.
  \end{itemize}
  \item Consequence: for a linear code the \textbf{minimum distance equals the minimum weight} of
        any non-zero codeword, $d' = w'$.
        \mk{This makes $d'$ far cheaper to compute: check $M$ weights instead of $M^2$ distances.}
  \item \textbf{Careful}: the presence of an all-zero codeword is \textbf{necessary but not
        sufficient} for linearity.
  \item \textbf{Matrix description}: a \textbf{generator matrix} $G$ ($k \times n$) encodes an
        information word $i$ of length $k$ into a codeword $c$ of length $n$:
        \[ c = i\,G \]
\end{itemize}

\lead{Cyclic codes}
\begin{itemize}
  \item A code $C$ is \textbf{cyclic} if it is \textbf{linear} \emph{and} \textbf{any cyclic shift}
        of a codeword is also a codeword.
  \item Ie if $a_0 a_1 \dots a_{n-1} \in C$ then $a_{n-1} a_0 \dots a_{n-2} \in C$.
  \item Eg $C_1 = \{0000, 0101, 1010, 1111\}$ is cyclic.
        $C_2 = \{0000, 0110, 1001, 1111\}$ is \textbf{not}, though it is
        \textbf{equivalent} to $C_1$: swapping the third and fourth components of $C_2$ gives
        $C_1$.
  \item \textbf{($+$)} Encoding and syndrome computation reduce to
        \mk{polynomial division, which is a simple shift register}. That is why nearly every
        practical block code is cyclic.
\end{itemize}}

\Q{Explain the 7-bit Hamming code method with a relevant example}{\m{5}}
{\color{sub}\footnotesize \tP{1} \yr{\textbf{80 Ch}} \gap$\cdot$\gap
\mk{worked both ways (encode and detect) in the Ch 6 companion}}

\sbsr{0.52}{%
\begin{itemize}
  \item A \textbf{linear block code} able to \textbf{identify and correct any single-bit error}
        within a block. Written as an $(n,k)$ Hamming code.
  \item Uses \textbf{modulo-2 arithmetic}, ie XOR.
  \item Parameters, for any integer $m \ge 3$:
  \begin{itemize}
    \item block length \; $n = 2^m - 1$
    \item message bits \; $k = 2^m - m - 1$
    \item parity check bits \; $n - k = m$
    \item minimum distance \; $d_{\min} = 3$, so error-correcting capability
          $t = \lfloor (d_{\min}-1)/2 \rfloor = 1$
    \item code efficiency (rate) \; $= k/n$
  \end{itemize}
  \item For $m = 3$: $n = 7$, $k = 4$. \mk{That is the ``7-bit Hamming code'' the paper asks for.}
  \item \textbf{($-$)} Offers \mk{little protection against burst errors}, since it corrects only
        one bit per block. Interleaving (Ch 5) is the standard fix.
\end{itemize}

\lead{Structure}
\begin{itemize}
  \item Parity bits are \textbf{inserted in between} the data bits, at positions
        $1, 2, 4, \dots$ (the powers of 2).
  \item Layout, bit 1 on the right:
        \; $D_7\,D_6\,D_5\,P_4\,D_3\,P_2\,P_1$
  \item Each parity bit covers the positions whose binary index includes its own bit:
  \begin{itemize}
    \item $P_1$ covers positions \textbf{1, 3, 5, 7}
    \item $P_2$ covers positions \textbf{2, 3, 6, 7}
    \item $P_4$ covers positions \textbf{4, 5, 6, 7}
  \end{itemize}
\end{itemize}}{%
\figT{c6_hamming_struct.png}
\vspace{2pt}

\lead{Worked example 1: encode 1011 with even parity}
\begin{itemize}
  \item Data goes into positions 7, 6, 5, 3: \; $D_7{=}1,\;D_6{=}0,\;D_5{=}1,\;D_3{=}1$.
  \item $P_1 = $ parity of positions 1,3,5,7 $= (P_1,1,1,1)$. Three ones already, so
        \textbf{$P_1 = 1$} to make it even.
  \item $P_2 = $ parity of positions 2,3,6,7 $= (P_2,1,0,1)$. Two ones, so \textbf{$P_2 = 0$}.
  \item $P_4 = $ parity of positions 4,5,6,7 $= (P_4,1,0,1)$. Two ones, so \textbf{$P_4 = 0$}.
  \item Transmitted codeword (bit 7 first): \; \mk{\textbf{1 0 1 0 1 0 1}}
\end{itemize}

\lead{Worked example 2: check the received word 1011011}
\begin{itemize}
  \item $P_1$ group (positions 1,3,5,7) reads \textbf{1011} $\rightarrow$ \textbf{odd} parity
        $\rightarrow$ error, set $P_1 = 1$.
  \item $P_2$ group (positions 2,3,6,7) reads \textbf{1001} $\rightarrow$ \textbf{even} parity
        $\rightarrow$ no error, set $P_2 = 0$.
  \item $P_4$ group (positions 4,5,6,7) reads \textbf{1101} $\rightarrow$ \textbf{odd} parity
        $\rightarrow$ error, set $P_4 = 1$.
  \item \textbf{Syndrome} $= P_4 P_2 P_1 = 101 = 5$ in decimal.
  \item \mk{So the received word is incorrect, and bit position 5 is the one in error.} Flip it to
        correct the word.
  \item \textbf{That is the whole elegance of the Hamming code}: the syndrome read as a binary
        number \textbf{is} the position of the faulty bit.
\end{itemize}}

\creamq{Other block codes}{\pill{gdL}{gd}{syllabus, no PYQ yet}}

\sbs{%
\begin{itemize}
  \item \textbf{Hadamard codes}
  \begin{itemize}
    \item Codewords are the \textbf{rows of a Hadamard matrix} $A$: an $N \times N$ matrix of 1s
          and 0s in which every row differs from every other in \textbf{exactly $N/2$} places.
    \item One row is all zeros; the rest have $N/2$ zeros and $N/2$ ones.
    \item Minimum distance $d_{\min} = N/2$, so it is a \mk{very strong but very low rate} code.
    \item Used for the \textbf{Walsh functions} that separate IS-95 CDMA channels.
  \end{itemize}
  \item \textbf{Golay codes}
  \begin{itemize}
    \item Linear binary \textbf{(23,12)} codes with $d_{\min} = 7$, so they correct
          \textbf{3 bits}.
    \item \mk{The only non-trivial perfect code that exists.} Every possible received word lies
          within distance 3 of exactly one codeword, so maximum likelihood decoding is exact.
  \end{itemize}
\end{itemize}}{%
\begin{itemize}
  \item \textbf{BCH (Bose--Chaudhuri--Hocquenghem) codes}
  \begin{itemize}
    \item \mk{One of the most powerful known classes of linear cyclic block codes.}
    \item A \textbf{subset of cyclic codes} whose generator polynomial roots are chosen carefully
          to give good error-correcting capability.
    \item Each BCH code corrects up to \textbf{$t$ random errors} per codeword.
    \item Parameters: block length $n = 2^m - 1$ ($m \ge 3$), message bits $k \ge n - mt$,
          minimum distance $d_{\min} \ge 2t+1$.
  \end{itemize}
  \item \textbf{Reed--Solomon (RS) codes}
  \begin{itemize}
    \item A subclass of \textbf{non-binary} BCH codes.
    \item \mk{The encoder operates on multi-bit symbols, not individual bits.}
    \item For $m$-bit symbols: block length $n = 2^m - 1$ symbols, message $k$ symbols, parity
          $n-k = 2t$ symbols, $d_{\min} = 2t+1$.
    \item \textbf{($+$)} Because a whole symbol counts as one error however many of its bits are
          wrong, RS codes are \mk{excellent against burst errors}.
    \item Used in CD / DVD, deep space links, and digital broadcasting.
  \end{itemize}
\end{itemize}}

\Q{Explain convolutional encoding and decoding, and the Viterbi algorithm}{\m{5}\m{3}}
{\color{sub}\footnotesize \tF{5} \yr{\textbf{74 Bh}, \textbf{75 Bh}, 70 Ma, \textbf{72 Ash},
\textbf{73 Bh}} \gap$\cdot$\gap
\mk{Viterbi alone is asked in 4 papers} $\rightarrow$ worked example in the Ch 6 companion}

\sbsr{0.53}{%
\lead{How convolutional codes differ from block codes}
\begin{itemize}
  \item In a block code, each codeword depends \textbf{only} on the current block of $k$ message
        digits.
  \item In a convolutional code, the block of \textbf{$n$ code digits} produced at each time unit
        depends on the current $k$ message digits \textbf{and on the previous ones held in
        memory}.
  \item \mk{So the encoder has memory. That is the whole difference.}
  \item No fixed block boundary is needed, which suits a \textbf{continuous} data stream.
\end{itemize}

\lead{Encoder parameters +Fig}
\begin{itemize}
  \item \textbf{Constraint length $K$} $=$ number of symbols the shift register can store,
        ie how far back the memory reaches. For the figure, $K = 3$.
  \item If the register stores $m$ previous information frames of length $k_0$, then
        \textbf{$v = m k_0$}.
  \item $k = 1$ message bit in, $n = 2$ encoded bits out.
  \item \textbf{Code rate} $= k/n = 1/2$, \textbf{code dimension} $= (n,k) = (2,1)$.
  \item Outputs are \textbf{mod-2 sums (XOR)} of selected register stages:
        \[ X_1 = M + M_1 + M_2, \qquad X_2 = M + M_2 \]
  \item \mk{The tap pattern is the code.} Two encoders with the same $K$ but different taps are
        different codes with different performance.
\end{itemize}

\lead{Three ways to represent the same encoder}
\begin{itemize}
  \item \textbf{Code tree}: every branch is one input symbol, labelled with the output pair.
        \textbf{Input 0 takes the upper branch, input 1 the lower.}
        \mk{($-$) Grows exponentially}, so it is unusable beyond a few bits.
  \item \textbf{State diagram}: nodes are the register contents, arcs are transitions. Compact,
        but hides time.
  \item \textbf{Trellis}: the state diagram \textbf{drawn against time}.
        \mk{A compact representation of the code tree}: the tree's repeating structure is folded
        back onto the same set of states at every step. This is what the decoder actually walks.
\end{itemize}}{%
\figT{c6_conv_enc.png}
\vspace{1pt}
{\color{sub}\footnotesize Rate 1/2, $K=3$ convolutional encoder. Two mod-2 adders tap the shift
register.}
\vspace{3pt}
\figT{c6_code_tree.png}
\vspace{1pt}
{\color{sub}\footnotesize Code tree for the message 110. States $a{=}00$, $b{=}01$, $c{=}10$,
$d{=}11$.}
\vspace{3pt}
\figT{c6_trellis.png}
\vspace{1pt}
{\color{sub}\footnotesize Code trellis and the equivalent state diagram. Solid line $=$ input 0,
dashed $=$ input 1.}}

\creamq{Viterbi decoding}{\m{3}\m{5}}
\sbsr{0.53}{%
\begin{itemize}
  \item Used for decoding convolutional codes. Also known as
        \textbf{maximum likelihood decoding}, because it finds the
        \mk{single most likely transmitted sequence}, not the most likely symbol.
  \item It is the practical implementation of \textbf{MLSE} (Ch 5).
\end{itemize}

\lead{The algorithm}
\begin{itemize}
  \item Walk the \textbf{trellis} one time step at a time.
  \item At each step, every node has \textbf{two incoming branches} (for a rate 1/2 code).
  \item For each branch compute a \textbf{metric}: the \textbf{Hamming distance} between the
        branch's expected output and the actually received pair. (Soft-decision decoders use
        Euclidean distance instead.)
  \item Add the branch metric to the surviving path metric that led into it.
  \item \mk{Keep only the smaller of the two, and discard the other.} The kept one is the
        \textbf{survivor} for that node.
  \item Repeat for every node at every step. \textbf{The number of survivors never grows}, which
        is exactly why the algorithm is tractable while an exhaustive tree search is not.
  \item At the end, trace back the survivor with the lowest total metric. Its branch labels give
        the decoded message.
\end{itemize}

\lead{Why it is optimal}
\begin{itemize}
  \item Discarding the worse of two paths into a node is safe: \mk{any future extension of the
        worse path can be applied equally to the better one}, so the worse path can never win
        later.
  \item Complexity is \textbf{linear in the message length} but \textbf{exponential in the
        constraint length $K$}, which is why practical $K$ stays around 7 to 9.
\end{itemize}}{%
\figT{c6_viterbi.png}
\vspace{1pt}
{\color{sub}\footnotesize Decoding $Y = 11\;01\;10$. Path metrics accumulate along the trellis;
at each node the larger metric is discarded.}}

\Q{Define turbo coding}{\m{2}\m{5}}
{\color{sub}\footnotesize \tP{2} \yr{\textbf{\texttt{82 Bh}}, 74 Ma}}

\sbsr{0.55}{%
\begin{itemize}
  \item A turbo code \textbf{combines the capability of convolutional codes with channel
        estimation theory}. It can be thought of as \textbf{nested or parallel convolutional
        codes}.
  \item \textbf{Structure} +Fig: two \textbf{recursive systematic convolutional} encoders
        $C_1$ and $C_2$, working on the \textbf{same} input bits but in \textbf{different
        orders}.
  \begin{itemize}
    \item $M$ in the figure is a \textbf{memory register}.
    \item A \textbf{delay line} and an \textbf{interleaver} (which reorders the information bits)
          force the input $d_k$ to appear in different sequences at the two encoders.
    \item \mk{The interleaver is what makes it a turbo code.} Without it the two encoders would
          see the same bit pattern and fail on the same errors.
    \item Because the encoders are \textbf{systematic}, the input sequence $d_k$ appears
          \textbf{directly at the output} $x_k$ alongside the parity streams $y_{1k}$ and
          $y_{2k}$.
  \end{itemize}
  \item \textbf{Decoding} is \textbf{iterative}: two soft-input soft-output decoders exchange
        \textbf{extrinsic information} and refine each other's estimates over
        $n_1$ and $n_2$ iterations. The name comes from this feedback, by analogy with a
        turbocharger. \added
  \item \textbf{Why it matters}: coding gains \textbf{far superior} to all previous error
        correcting codes, letting a wireless link come
        \mk{amazingly close to the Shannon capacity bound}.
  \item \textbf{($-$)} \textbf{Large decoding delay} from the interleaver and the iterations, so
        it suits data far better than real-time voice.
\end{itemize}

\lead{Coding gain \pill{gdL}{gd}{syllabus, no PYQ yet}}
\begin{itemize}
  \item For a given modulation and target BER, \textbf{coding gain} $=$ the ratio of the
        \textbf{SNR per bit required uncoded} to the \textbf{SNR per bit required coded}:
        \[ G(\text{dB}) = \left(\frac{E_b}{N_0}\right)_{\text{uncoded}}(\text{dB})
                        - \left(\frac{E_b}{N_0}\right)_{\text{coded}}(\text{dB}) \]
  \item Ie \mk{how much transmit power the code saves you}.
  \item Every error control code has its own coding gain, depending on the code, the decoder
        implementation and the channel BER.
\end{itemize}

\lead{Trellis coded modulation (TCM) \pill{gdL}{gd}{syllabus, no PYQ yet}}
\begin{itemize}
  \item Traditionally coding and modulation are designed \textbf{separately}, even though both
        exist to fight channel errors.
  \item TCM \textbf{integrates the channel encoder with the modulator}, so coding gain is obtained
        \mk{without any bandwidth expansion}.
  \item Uses \textbf{signal set expansion} (eg 8-PSK instead of 4-PSK) to supply the redundancy,
        then designs the coding and the signal mapping \textbf{jointly} to maximise the
        \textbf{free distance}, ie the minimum Euclidean distance between coded sequences.
  \item Decoded by a \textbf{soft decision maximum likelihood sequence decoder}, ie Viterbi again.
  \item Gains as large as \textbf{6 dB} with \textbf{no bandwidth expansion and no reduction in
        information rate}.
\end{itemize}}{%
\figT{c6_turbo.png}
\vspace{1pt}
{\color{sub}\footnotesize Turbo encoder. Two systematic convolutional encoders $C_1$ and $C_2$;
the interleaver ahead of $C_2$ makes the two see different bit orders.}}
